% Options for packages loaded elsewhere
\PassOptionsToPackage{unicode}{hyperref}
\PassOptionsToPackage{hyphens}{url}
\PassOptionsToPackage{dvipsnames,svgnames,x11names}{xcolor}
\documentclass[
  ngerman,
  10pt,
  a4paper,
]{article}
\usepackage{xcolor}
\usepackage[margin=22mm]{geometry}
\usepackage{amsmath,amssymb}
\setcounter{secnumdepth}{-\maxdimen} % remove section numbering
\usepackage{iftex}
\ifPDFTeX
  \usepackage[T1]{fontenc}
  \usepackage[utf8]{inputenc}
  \usepackage{textcomp} % provide euro and other symbols
\else % if luatex or xetex
  \usepackage{unicode-math} % this also loads fontspec
  \defaultfontfeatures{Scale=MatchLowercase}
  \defaultfontfeatures[\rmfamily]{Ligatures=TeX,Scale=1}
\fi
\usepackage{lmodern}
\ifPDFTeX\else
  % xetex/luatex font selection
  \setmainfont[Path=/Users/stefanhamann/.cache/codex-runtimes/codex-primary-runtime/dependencies/native/libreoffice-headless/libreoffice/LibreOfficeDev.app/Contents/Resources/fonts/truetype/,Extension=.ttf,BoldFont=DejaVuSans-Bold,ItalicFont=DejaVuSans-Oblique,BoldItalicFont=DejaVuSans-BoldOblique]{DejaVuSans}
  \setmonofont[Path=/Users/stefanhamann/.cache/codex-runtimes/codex-primary-runtime/dependencies/native/libreoffice-headless/libreoffice/LibreOfficeDev.app/Contents/Resources/fonts/truetype/,Extension=.ttf,BoldFont=DejaVuSansMono-Bold,ItalicFont=DejaVuSansMono-Oblique,BoldItalicFont=DejaVuSansMono-BoldOblique]{DejaVuSansMono}
\fi
% Use upquote if available, for straight quotes in verbatim environments
\IfFileExists{upquote.sty}{\usepackage{upquote}}{}
\IfFileExists{microtype.sty}{% use microtype if available
  \usepackage[]{microtype}
  \UseMicrotypeSet[protrusion]{basicmath} % disable protrusion for tt fonts
}{}
\makeatletter
\@ifundefined{KOMAClassName}{% if non-KOMA class
  \IfFileExists{parskip.sty}{%
    \usepackage{parskip}
  }{% else
    \setlength{\parindent}{0pt}
    \setlength{\parskip}{6pt plus 2pt minus 1pt}}
}{% if KOMA class
  \KOMAoptions{parskip=half}}
\makeatother
\ifLuaTeX
\usepackage[bidi=basic,shorthands=off]{babel}
\else
\usepackage[bidi=default,shorthands=off]{babel}
\fi
\ifPDFTeX
\else
\babelfont{rm}[Path=/Users/stefanhamann/.cache/codex-runtimes/codex-primary-runtime/dependencies/native/libreoffice-headless/libreoffice/LibreOfficeDev.app/Contents/Resources/fonts/truetype/,Extension=.ttf,BoldFont=DejaVuSans-Bold,ItalicFont=DejaVuSans-Oblique,BoldItalicFont=DejaVuSans-BoldOblique]{DejaVuSans}
\fi
\ifLuaTeX
  \usepackage{selnolig} % disable illegal ligatures
\fi
\setlength{\emergencystretch}{3em} % prevent overfull lines
\providecommand{\tightlist}{%
  \setlength{\itemsep}{0pt}\setlength{\parskip}{0pt}}
\usepackage{fvextra}
\RecustomVerbatimEnvironment{verbatim}{Verbatim}{breaklines=true,fontsize=\footnotesize}
\usepackage{xurl}
\usepackage{seqsplit}
\usepackage{adjustbox}
\usepackage{ragged2e}
\usepackage{titlesec}
\titleformat{\section}{\large\bfseries\raggedright}{}{0em}{}
\titleformat{\subsection}{\normalsize\bfseries\raggedright}{}{0em}{}
\DeclareRobustCommand{\codeword}[1]{\texttt{\seqsplit{#1}}}
\AtBeginEnvironment{longtable}{\small\RaggedRight}
\usepackage{newunicodechar}
\newunicodechar{𝒞}{\ensuremath{\mathcal C}}
\newunicodechar{𝔊}{\ensuremath{\mathfrak G}}
\newunicodechar{𝔤}{\ensuremath{\mathfrak g}}
\newunicodechar{𝓗}{\ensuremath{\mathcal H}}
\newunicodechar{ℋ}{\ensuremath{\mathcal H}}
\newunicodechar{𝔄}{\ensuremath{\mathfrak A}}
\newunicodechar{𝟙}{\ensuremath{\mathbf 1}}
\newunicodechar{𝕀}{\ensuremath{I}}
\newunicodechar{𝕋}{\ensuremath{\mathbb T}}
\newunicodechar{𝟘}{\ensuremath{\mathbf 0}}
\newunicodechar{⊞}{\ensuremath{\boxplus}}
\newunicodechar{∘}{\ensuremath{\circ}}
\newunicodechar{⟦}{\ensuremath{[\![}}
\newunicodechar{⟧}{\ensuremath{]\!]}}
\setlength{\emergencystretch}{4em}
\setlength{\parindent}{0pt}
\setlength{\parskip}{0.45em}
\AtBeginDocument{\hypersetup{colorlinks=true,linkcolor=black,urlcolor=blue}\pdfstringdefDisableCommands{\def\codeword#1{#1}}\RaggedRight}
\usepackage{fancyhdr}
\pagestyle{fancy}
\fancyhf{}
\fancyhead[L]{\small TFPT / Universalraum - v1.6.7}
\fancyfoot[R]{\thepage}
\setlength{\headheight}{15pt}
\usepackage{bookmark}
\IfFileExists{xurl.sty}{\usepackage{xurl}}{} % add URL line breaks if available
\urlstyle{same}
\hypersetup{
  pdflang={de-DE},
  colorlinks=true,
  linkcolor={Maroon},
  filecolor={Maroon},
  citecolor={Blue},
  urlcolor={Blue},
  pdfcreator={LaTeX via pandoc}}

\author{}
\date{}

\begin{document}

\section{TFPT / Universalraum - Update
v1.6.7}\label{tfpt-universalraum---update-v1.6.7}

\begin{enumerate}
\def\labelenumi{\arabic{enumi}.}
\setcounter{enumi}{14}
\tightlist
\item
  September 2026. Ausschließlich Änderungen gegenüber v1.6.6.
\end{enumerate}

\subsection{Neue Priorität: derselbe Prozess hinter den
Schatten}\label{neue-priorituxe4t-derselbe-prozess-hinter-den-schatten}

Die drei zusätzlich eingesandten Grundsatztexte und die Holografiefrage
verschieben die Priorität: Die gemeinsame Herkunft von Operation,
Zustand und Auslese ist wichtiger als weitere Dezimalstellen einer
lokalen Rechnung. Alle sechs Quellen dieser Runde stehen unverändert im
Hauptdokument; die neuen Herleitungen und Grenzen stehen in dessen Teil
B.

\textbf{Ein nativer Kodierer existiert bereits.} Aus WW dagger gleich 8
I60 folgt die Isometrie V gleich W dagger durch Wurzel 8 vom 60er-Träger
in den 2016-dimensionalen Fermionpaarraum. Der dunkle Kern der
Gegenrichtung hat Dimension 1956. Im hellen N=2-Sektor ist die
gemeinsame Dynamik eine Zweiermatrix mit Kopplung Wurzel 8 mal g,
jeweils für denselben inneren Typ. Dies ist eine interne Kodierung, noch
keine holografische Raumzeit.

Der Code schützt nicht automatisch gegen den Verlust einer einzelnen
Mode: Für jede der 64 Moden hat V dagger n\_r V die Eigenwerte 0
(45-mal) und 1/8 (15-mal), ist also nicht skalar. Eine verlorene Mode
kann logische Information verraten. Auch viele Schatten bestimmen nicht
immer den Ursprung: Zwei orthogonale GHZ-Zustände besitzen identische
sämtliche echten Teilmarginalen. Ergänzende kalibrierte Qubitansichten
können dagegen gemeinsam vollständig sein. Entscheidend ist der
gemeinsam unsichtbare Kern.

\subsection{Neue exakte Prüfungen der einfachsten
Dynamik}\label{neue-exakte-pruxfcfungen-der-einfachsten-dynamik}

\begin{itemize}
\tightlist
\item
  Multiplizitäten sind mögliche Dynamikträger, aber noch keine Orte.
  Schon N=3 besitzt zwei symmetriegleiche Zusammensetzungen mit echter
  interner Umwandlung. Der volle G-Kommutant hat dort Dimension 16,
  nicht die sehr große Dimension des Kommutanten eingeschränkter
  Kontrollen.
\item
  Zeitabhängige Kreuzkorrelation ist kein hinreichender
  Transportnachweis. Ein stationärer verschränkter Zustand zweier
  ungekoppelter Qubits liefert C\_AB(t)=sin(t), also Anfangswert null,
  ohne irgendeine lokale Signalwirkung. Gefordert wird eine Änderung der
  unbedingten B-Statistik nach A-Eingriff.
\item
  Der native 480-Kanten-Paargraph ist zusammenhängend. Bei
  festgehaltenen Bosonen bleiben nur triviale und globale
  Fermionteilmengenparität.
\item
  Derselbe W und seine reelle invariante Bosonpaarung erlauben einen
  quadratischen Bosonpaarterm sowie einen kubischen
  Boson-plus-Fermionpaar- Erzeugungsterm. Beide ändern N um vier und
  erhalten G und dessen Z4-Zentrum. Die Zahl kontinuierlicher diagonaler
  Phasensymmetrien fällt exakt von neun auf acht. Zulässigkeit ist keine
  native Erzeugbarkeit.
\item
  Ein ganzzahliger konstanter Ladungslift der vollständigen genannten
  Z4-Klammern ist unmöglich. Eine wörtliche Identifikation der
  Grad-eins- Lieklammer mit dem CAR-Kommutator primitiver Fermionen
  scheitert ebenfalls. Das schließt andere zusammengesetzte
  Realisierungen nicht aus.
\end{itemize}

Auf der reellen phasengleichen Hamiltonfamilie lassen sich Bosonpaarterm
kappa und zweiter kubischer Kanal lambda durch dieselbe uniforme
Bogoliubov-Transformation genau dann beide entfernen, wenn g Quadrat
größer lambda Quadrat ist und

\[\adjustbox{max width=\linewidth}{$\displaystyle 
\kappa(g^2+\lambda^2)-2\Delta g\lambda=0.
$}\]

Dann kann eine andere Teilchenzahl erhalten sein, obwohl die
ursprüngliche nicht erhalten ist. Echte neue Dynamik muss von dieser
Variablenfreiheit getrennt werden. Selbst bei echter N-Brechung ist mu N
weiterhin symmetrieverträglich: Z4 allein wählt das Vakuum nicht aus.
Die bisherige Grundzustands- und Polzertifizierung wird auf keine
Erweiterung übertragen.

\subsection{Vier statt einer Richtung - jetzt
konstruiert}\label{vier-statt-einer-richtung---jetzt-konstruiert}

Der vollständige Singulettbereich im N=64-Sektor bei zwei Bosonen
besitzt exakt vier Dimensionen. Der fertige neue Worker-Bericht und
unsere unabhängige Charakterrechnung stimmen überein. Die Typen auf der
Bosonseite sind (54,1), (1,20'), (54,20') und (45,15). Durch Projektion
des frisch berechneten nativen Zustands v2 erhält man alle vier
Basisrichtungen.

Ihre quadrierten Kopplungen vom Ein-Boson-Singulett betragen
\(g^2(36,16,504,360)\). Die Summe ist \(916g^2\). Der ganze Bereich bis
Bosonstufe zwei ist damit ein exakt bestimmter Sechszustands-Ausschnitt:
eine helle Dreierkette und drei nach unten dunkle Richtungen. Höhere
Bosonstufen bleiben angekoppelt; der Ausschnitt ist keine vollständige
Lösung.

\subsection{Die richtige einfache Kette und eine neue
Energiegrenze}\label{die-richtige-einfache-kette-und-eine-neue-energiegrenze}

Die Lanczos-Kette der Wiederholungen des ganzen H vermeidet die falsche
Annahme, ein Kettenschritt sei identisch mit einer Bosonzahl. Neu
bestimmt:

\[\adjustbox{max width=\linewidth}{$\displaystyle 
b_4^2/g^2=78877653/47632,\qquad
a_4/\Delta=105168998/26292551.
$}\]

Der vierte neue Vektor enthält sowohl die Stufe vier als auch den
bekannten Seitenzweig auf Stufe zwei. Die ersten fünf Kettenglieder
bestimmen die zehn Energiemomente \(\langle F,H^nF\rangle\), n von 0 bis
9, exakt. Dies sind Momente der gefüllten Referenz, nicht die geladenen
Momente des Grundzustands.

Eine rationale Eigenwerteinschließung liefert \(E_0<-1.12963811\Delta\).
Mit dem früheren N=63-Floor folgt \(\epsilon>0.00773911\Delta\). Das
verbessert die frühere untere Polgrenze 0.007737 geringfügig; Polgewicht
und obere Polgrenze werden dadurch nicht verändert. Die vier-Boson-Norm
bleibt ein früherer exakter Eingang, nicht eine in dieser Runde frisch
wiederholte Enumeration.

\subsection{Der bilineare Rest ist vollständig
zerlegt}\label{der-bilineare-rest-ist-vollstuxe4ndig-zerlegt}

\[\adjustbox{max width=\linewidth}{$\displaystyle 
\operatorname{End}(16\otimes4)=
(1,1)\oplus(45,1)\oplus(210,1)
\oplus(1,15)\oplus(45,15)\oplus(210,15).
$}\]

Die 4035 zuvor unbestimmten Komponenten sind 210 plus 675 plus 3150. Die
Zerlegung ist innere Darstellungstheorie, keine neue Teilchenzählung.
Sie wurde durch Clifford-Gram und Weyl-Multiplizitäten abgesichert.

\subsection{Korrekturen und
Reichweite}\label{korrekturen-und-reichweite}

\begin{itemize}
\tightlist
\item
  Der Clock liegt in Spin(10). Nichtskalare Wirkung widerlegt diese
  Zugehörigkeit nicht; die gegenteilige Eingabe wird korrigiert.
\item
  Die zusätzlichen SU(4)- und Spin(10)-Stufen haben im N=3-Sektor die
  Kommutantdimensionen 2247168 und 1648. Verfügbarkeit ist nicht binär.
\item
  Volle Gruppe plus alle Besetzungen erzeugt bedingt die volle
  Sektoralgebra. Ohne gewährte Gruppenoperationen bleibt trotz
  Besetzungsauslese eine nichtskalare Cartan-Erhaltung. Eine Auslese ist
  noch kein Transferinstrument.
\item
  Drei Feldkanäle mit Normquadrat 64 ergeben Normquadrat 192 und Norm
  \(8\sqrt3\), nicht 24. Ein Projektorranganteil von einem Drittel ist
  kein universelles Spektralgewicht.
\item
  Der fertig eingefrorene Grundzustandsbericht meldet 379 exakte und
  drei numerische Guards. Seine schwächeren Sektorschranken widerlegen
  den stärkeren früheren Grundzustandssatz nicht. Beim Operationsbericht
  gibt es eine Checkerhash-Abweichung und 306 statt angegebener 307
  Guards; seine entscheidenden Aussagen wurden deshalb unabhängig
  kontrolliert.
\end{itemize}

\subsection{Fundamentale Auswahlfrage jetzt als konkreter
Test}\label{fundamentale-auswahlfrage-jetzt-als-konkreter-test}

Bei gleichem W und denselben Symmetrien ergibt sich für dieselbe
Referenz

\[\adjustbox{max width=\linewidth}{$\displaystyle 
\mu_4\mu_2/\mu_3^2=1+1396(g/\Delta)^2.
$}\]

Die Werte 4,49 und 1,8725 für die beiden zulässigen Prüfparameter 1/20
und 1/40 zeigen: Die bisherigen geometrischen Grunddaten wählen noch
keine eindeutige Dynamik. Eine zusätzliche einfache Ursprungsregel
könnte das leisten; sie ist bislang nicht konstruiert.

\textbf{Prüfung:} 822 neue exakte Prüfbedingungen je Python-Modus,
identische Ergebnisdateien, frisch wiederholte historische
C++-Kontraktion, 24 eingefrorene Eingaben. Kein vollständiger
Fremdprogramm-Replay, kein geschlossenes T1-T8-Tor, keine fertige TOE.

\textbf{Nächste Schritte:} den minimalen Ursprungsvertrag auswählen; die
gemeinsame Schattenkarte samt verbleibender Informationslücken aufbauen;
darauf einen kausalen Eingriffstransfer mit demselben relativistischen
Feldadapter nachweisen. Die H-Kette bleibt ein kontrolliertes
Hilfswerkzeug. Vollständige Herleitungen und frühere Ergebnisse stehen
im Hauptdokument v1.6.7.

\end{document}
